\documentclass{article}
\usepackage{graphicx}
\usepackage{xcolor}
\usepackage{amsmath}
\usepackage{amssymb}

\setlength{\fboxsep}{4pt}
\setlength{\parindent}{0pt}

\title{Linear Algebra Done Right Solutions}
\author{Vadym}
\date{June 2026}

\begin{document}

\maketitle
\thispagestyle{empty}
\clearpage
\setcounter{page}{1}

\section*{Exercises 1C}
{\color{red}\rule{\textwidth}{0.5pt}}

\subsection*{2. Verify all assertions about subspaces in Example 1.35.}

\subsubsection*{(a)}

\noindent\fbox{
\parbox{\textwidth}{
If $b \in F$, then
\[
\{(x_1,x_2,x_3,x_4) \in F^4 : x_3 = 5x_4 + b\}
\]
is a subspace of $F^4$ if and only if $b = 0$.
}
}

\subsubsection*{\fbox{Proof}}
\subsubsection*{Forward direction}
\noindent
Suppose that
\[
U=\{(x_1,x_2,x_3,x_4)\in F^4 : x_3=5x_4+b\}
\]
is a subspace of $F^4$.

Since $U$ is a subspace of $F^4$, it contains the zero vector of $F^4$. Hence
\[
(0,0,0,0)\in U.
\]
Since every vector in $U$ satisfies
\[
x_3 = 5x_4 + b,
\]
the zero vector must satisfy this equation as well. Thus
\[
\begin{aligned}
0 &= 5(0) + b\\
0 &= b,
\end{aligned}
\]
and hence $b=0$.

\subsubsection*{Reverse Direction}

\noindent
Suppose that $b \neq 0$. We show that $U$ is not a subspace of $F^4$.

Assume, for the sake of contradiction, that $U$ is a subspace of $F^4$. Since $U$ is a subspace of $F^4$, it contains the zero vector of $F^4$. Hence

\[
(0,0,0,0)\in U.
\]

Since every vector in $U$ satisfies

\[
x_3 = 5x_4 + b,
\]

the zero vector must satisfy this equation as well. Thus

\[
\begin{aligned}
0 &= 5(0) + b \\
0 &= b.
\end{aligned}
\]

This contradicts the assumption that $b \neq 0$. Therefore, $U$ is not a subspace of $F^4$.

\hfill$\square$
\end{document}